\documentclass[12pt,a4paper]{article}

\usepackage{amsmath,amssymb}
\usepackage{graphicx}
\usepackage{booktabs}
\usepackage{array}
\usepackage{hyperref}
\usepackage{natbib}
\usepackage{geometry}
\usepackage{xcolor}
\usepackage{microtype}
\usepackage{multirow}
\usepackage{float}
\usepackage{colortbl}
\usepackage{longtable}

\geometry{margin=1in}

\definecolor{mcmcrow}{RGB}{220,230,250}
\definecolor{newresult}{RGB}{200,230,200}
\definecolor{weakrow}{RGB}{250,220,220}

\hypersetup{
  colorlinks=true,
  linkcolor=blue,
  citecolor=blue,
  urlcolor=blue
}

% ─────────────────────────────────────────────────────────────────────────────
\title{\textbf{The Thermodynamic Identity Governing the Virial Theorem:\\
Derivation from First Principles and Evidence\\
Across 37 Orders of Magnitude}}

\author{Heath W.\ Mahaffey\\
\small Independent Researcher, Entiat, Washington 98822, USA\\
\small \href{mailto:hmahaffeyges@gmail.com}{hmahaffeyges@gmail.com}\\
\small \url{https://github.com/hmahaffeyges/IAM-Validation}}

\date{March 2026\\
\small Zenodo DOI:
\href{https://doi.org/10.5281/zenodo.18702042}{10.5281/zenodo.18702042}}
% ─────────────────────────────────────────────────────────────────────────────

\begin{document}
\maketitle

% ─────────────────────────────────────────────────────────────────────────────
\begin{abstract}
The virial theorem, established by Clausius in 1870, states that for any
stable bound system under a $1/r$ potential, the time-averaged kinetic and
potential energies satisfy $2\langle K \rangle + \langle V \rangle = 0$.
The theorem has been verified across 37 orders of magnitude in physical
scale, but one question has persisted: what is $\langle K \rangle$?
The virial theorem identifies $\langle K \rangle$ as the average kinetic
energy of bound particles. It does not identify what physical process
produces the $1/2$ partition, nor what happens thermodynamically to the
energy released during the transition from the unbound to the bound state.

We show that the virial theorem and Landauer's principle are two
expressions of the same underlying thermodynamic identity. The first law,
the second law, Landauer's principle, and the equilibrium condition of the
$1/r$ potential together require:
\begin{equation*}
E_\mathrm{Landauer} = Q = -E_f = \langle K \rangle,
\end{equation*}
where $Q$ is the heat released irreversibly at the virialization boundary,
$\Delta S$ is the entropy produced, and $E_f$ is the total energy of the
bound state. This identifies $\langle K \rangle$ for the first time as a
thermodynamic quantity: the Landauer cost of the system's transition from
the unbound to the bound state. The thermodynamic identity is:
\begin{equation}
\langle K \rangle = Q = T\,\Delta S = E_\mathrm{Landauer}.
\label{eq:identity_abstract}
\end{equation}
The virial theorem $2\langle K \rangle + \langle V \rangle = 0$ is the
mechanical expression of this identity. The identity and the virial theorem
are two descriptions of the same physical condition at the boundary between
unbound and bound states.

The data are consistent with this identity across 37 orders of magnitude
in physical scale, from hydrogen atoms at $10^{-10}$~m to the
cosmological coupling constant $\beta_m = \Omega_m/2$, for which 17
independent MCMC chains against the full Planck~2018 likelihood
\citep{PlanckVI2020} return $\beta_m = 0.1583 \pm 0.0033$ against the
predicted $0.1577$, consistent at $0.2\sigma$ with zero free parameters.
The Informational Actualization Model (IAM)~\citep{IAM_theory} is one
framework that instantiates this identity at cosmological scales,
recovering $\beta_m = \Omega_m/2 = 0.1575$ with no fitted parameters.

Three of the four fundamental forces satisfy this identity within their
respective domains. The weak nuclear force is a structural exception,
excluded not by coincidence but by physical necessity: the force whose
function is to enable irreversible transitions between stable states cannot
simultaneously satisfy an equilibrium condition. This exception is itself
evidence for the identity's physical content.

We propose that this thermodynamic identity constitutes a fundamental law
of physics of which the virial theorem is the mechanical special case.
\end{abstract}

\noindent\textbf{Keywords:} virial theorem; Landauer's principle;
thermodynamic identity; gravitational decoherence; entropy; bound states;
cosmological constant; Informational Actualization Model

% ─────────────────────────────────────────────────────────────────────────────
\section{Introduction}
\label{sec:intro}
% ─────────────────────────────────────────────────────────────────────────────

\subsection{The unexplained 1/2}

The virial theorem has governed the energy partitioning of bound physical
systems since Clausius stated it mathematically in
1870~\citep{Clausius1870}. For any stable system bound by a $1/r$
potential, the time-averaged kinetic and potential energies satisfy:
\begin{equation}
2\langle K \rangle + \langle V \rangle = 0,
\qquad \langle K \rangle = \tfrac{1}{2}|\langle V \rangle|.
\label{eq:virial}
\end{equation}
The theorem is universally verified --- from the hydrogen atom
($10^{-10}$~m) to galaxy clusters ($10^{23}$~m) and, as shown below, to
the cosmological coupling constant ($10^{26}$~m) --- a span of 37 orders
of magnitude. It is one of the most broadly confirmed results in physics.

Yet a question has gone unanswered for 154 years: what is $K$?

The virial theorem identifies $K$ as the time-averaged kinetic energy of
the bound particles. It establishes the ratio $K/|V| = 1/2$ as a
mathematical consequence of the $1/r$ potential form via the equations of
motion. It does not explain why the energy balance at equilibrium requires
exactly this ratio, nor what physical process produces it, nor what happens
to the energy that is not retained as kinetic energy during the transition
from the unbound to the bound state.

Clausius derived Equation~\ref{eq:virial} from Newton's second law. The
derivation is mechanical. It describes the equilibrium state. It is silent
on the thermodynamic content of the transition that produces the
equilibrium state, and silent on the physical identity of the energy
released during that transition.

The thermodynamics of gravitational systems has a substantial history.
Lynden-Bell and Wood~\citep{LyndenBell1968} identified the gravothermal
catastrophe; Padmanabhan~\citep{Padmanabhan1990} analyzed statistical
mechanics of self-gravitating systems; Saslaw~\citep{Saslaw1985} developed
the statistical thermodynamics of gravitational clustering across stellar
and galactic scales. These works study the thermodynamic
properties of systems already in or near virial equilibrium. What has not
been shown is that the virial theorem and Landauer's principle are two
expressions of the same identity, and that this identity provides the
physical meaning of $\langle K \rangle$ as the Landauer cost of the
virialization transition itself.

\subsection{The present work}

This paper shows that the virial theorem is the mechanical expression of a
thermodynamic identity, and uses that identity to provide the first
physical identification of $\langle K \rangle$. The argument proceeds in
four steps:

\begin{enumerate}
\item The transition from an unbound to a bound state under a $1/r$
potential releases heat $Q$ irreversibly at the boundary. By the first law,
$Q = -E_f$.
\item By the second law and Landauer's principle~\citep{Landauer1961},
the thermodynamic cost of this irreversible transition is
$E_\mathrm{Landauer} = T_\mathrm{vir}\,\Delta S = Q$. The virial
temperature cancels exactly, making this result scale-invariant.
\item The equilibrium condition of the $1/r$ potential --- expressed by
the virial theorem itself --- requires $-E_f = \langle K \rangle$ at
stable bound equilibrium.
\item Combining steps 1--3: $E_\mathrm{Landauer} = Q = -E_f = \langle K \rangle$.
This identifies $\langle K \rangle$ as the Landauer cost of the
virialization transition. The virial theorem is the mechanical statement
of this thermodynamic identity.
\end{enumerate}

The result is that $K$ is not merely a mechanical quantity. It is the
thermodynamic cost of the system's transition from the unbound to the
bound state --- the heat released irreversibly and encoded permanently as
entropy at the virialization boundary. The virial theorem is the mechanical
expression of this thermodynamic identity.

Section~\ref{sec:derivation} presents the derivation.
Section~\ref{sec:statement} states the identity formally and compares it
to the virial theorem.
Section~\ref{sec:evidence} presents the evidence across 37 orders of
magnitude.
Section~\ref{sec:forces} analyzes the four fundamental forces.
Section~\ref{sec:mcmc} presents the MCMC confirmation as the strongest
quantitative test.
Section~\ref{sec:discussion} discusses implications.

% ─────────────────────────────────────────────────────────────────────────────
\section{Derivation}
\label{sec:derivation}
% ─────────────────────────────────────────────────────────────────────────────

\subsection{Setup}

Consider a system of total mass $M$ transitioning from an unbound state to
a stable bound equilibrium state under a $1/r$ gravitational potential. We
adopt the natural reference state for a $1/r$ potential: the unbound state
at rest at infinity, with total energy $E_i = 0$.

The bound equilibrium state has total energy:
\begin{equation}
E_f = \langle K \rangle + \langle V \rangle.
\label{eq:Ef}
\end{equation}
No assumption about the ratio $\langle K \rangle / |\langle V \rangle|$ is
made at this stage. The thermodynamic argument will identify the physical
content of that ratio once the equilibrium condition of the $1/r$ potential
is combined with Landauer's principle.

\subsection{Step 1: Heat released at the boundary (First Law)}

The transition from unbound to bound is irreversible. Gravitational
collapse is a one-way process; a virialized system does not spontaneously
return to the unbound state under the same potential. By the first law of
thermodynamics, the heat released irreversibly to the environment during
the transition is:
\begin{equation}
Q = E_i - E_f = -E_f = -(\langle K \rangle + \langle V \rangle).
\label{eq:Q_firstlaw}
\end{equation}
This is exact. No approximation has been made. $Q > 0$ requires
$E_f < 0$, which is satisfied for any stable bound state.

\subsection{Step 2: Entropy produced at the boundary (Second Law)}

The transition is irreversible. By the second law of thermodynamics, it
produces entropy:
\begin{equation}
\Delta S = \frac{Q}{T_\mathrm{vir}},
\label{eq:entropy}
\end{equation}
where $T_\mathrm{vir}$ is the virial temperature of the bound system. For
a self-gravitating system of $N$ particles with velocity dispersion
$\sigma$, the virial temperature is $T_\mathrm{vir} = m\sigma^2/3k_B$
from equipartition, where $m$ is the characteristic particle mass.

\subsection{Step 3: Landauer cost of the entropy production}

By Landauer's principle~\citep{Landauer1961}, the minimum thermodynamic
cost of producing entropy $\Delta S$ irreversibly is:
\begin{equation}
E_\mathrm{Landauer} = T_\mathrm{vir}\,\Delta S.
\label{eq:landauer}
\end{equation}
Substituting Equation~\ref{eq:entropy}:
\begin{equation}
E_\mathrm{Landauer} = T_\mathrm{vir} \times \frac{Q}{T_\mathrm{vir}} = Q.
\label{eq:landauer2}
\end{equation}
The virial temperature cancels exactly. The Landauer cost equals the heat
released at the boundary, independently of temperature. This is the first
key result: the thermodynamic cost of the transition is exactly $Q$,
regardless of the scale or temperature of the system.

\subsection{Step 4: The equilibrium condition of the $1/r$ potential}

Steps 1--3 establish that the Landauer cost of the virialization
transition equals the heat released at the boundary:
\begin{equation}
E_\mathrm{Landauer} = Q = -E_f = -(\langle K \rangle + \langle V \rangle).
\label{eq:combined}
\end{equation}
This is a real and nontrivial result: the thermodynamic cost of the
transition from the unbound to the bound state equals the binding energy
of the resulting system.

To complete the identification of $\langle K \rangle$, we invoke the
equilibrium condition of the $1/r$ potential. For a system governed by a
potential $V \propto r^{-1}$, Euler's theorem for homogeneous functions
of degree $-1$ requires:
\begin{equation}
2\langle K \rangle + \langle V \rangle = 0 \quad \Longrightarrow \quad
-E_f = -(\langle K \rangle + \langle V \rangle) = \langle K \rangle.
\label{eq:euler}
\end{equation}
This is the virial theorem stated as the equilibrium condition of the
$1/r$ potential --- not assumed as a result, but recognized as the
condition that defines what it means for a $1/r$ system to be at stable
bound equilibrium. Any system not satisfying Equation~\ref{eq:euler} is
either not yet at equilibrium or is unstable.

\subsection{Closing the argument}

Substituting Equation~\ref{eq:euler} into Equation~\ref{eq:combined}:
\begin{equation}
\boxed{E_\mathrm{Landauer} = Q = -E_f = \langle K \rangle = \tfrac{1}{2}|\langle V \rangle|}
\label{eq:identity}
\end{equation}
The Landauer cost of the virialization transition equals the heat released
at the boundary, which equals the binding energy of the system, which
equals the kinetic energy at equilibrium, which equals half the magnitude
of the potential energy. The chain is exact at every step.

This is not a derivation of the virial theorem from thermodynamics alone.
It is the demonstration that the virial theorem and Landauer's principle
are two expressions of the same identity: the virial theorem provides the
mechanical equilibrium condition ($-E_f = \langle K \rangle$); Landauer's
principle provides the thermodynamic identification of what that energy
physically is ($E_\mathrm{Landauer} = Q = -E_f$). Together they establish
that $\langle K \rangle$ is not merely a mechanical quantity but the
irreversible thermodynamic cost of the system's transition to the bound
state. The $1/2$ partition is the unique ratio at which mechanical and
thermodynamic equilibrium are simultaneously satisfied for a $1/r$
potential. Any other partition would be inconsistent with either the
virial equilibrium condition or the Landauer cost of the transition.

\subsection{The physical meaning of K}

Equation~\ref{eq:identity} identifies $\langle K \rangle$ for the first
time as a thermodynamic quantity. The kinetic energy of a virialized system
is not merely the average kinetic energy of its constituent particles. It
is equal in magnitude to the heat $Q$ released irreversibly at the
virialization boundary. $Q$ is encoded as entropy on the nearest available
surface at Landauer cost $k_B T \ln 2$ per bit. $\langle K \rangle$ is
retained in the system as the kinetic energy of the bound particles. Both
are exactly half of $|\langle V \rangle|$ --- one half leaving the system
as irreversible entropy, one half remaining as mechanical motion. These are
not two separate processes. They are the two faces of the same partition of
$|\langle V \rangle|$ at the virialization boundary.

% ─────────────────────────────────────────────────────────────────────────────
\section{The Thermodynamic Identity: Formal Statement}
\label{sec:statement}
% ─────────────────────────────────────────────────────────────────────────────

\subsection{The identity}

\begin{quote}
\textit{As a system crosses from an unbound to a bound state, it
immediately encounters a $1/r$ potential and releases heat $Q$
irreversibly. This heat equals the kinetic energy at equilibrium,
$Q = \langle K \rangle$, and $Q$ is encoded onto the nearest available
surface as entropy at thermodynamic cost $k_B T \ln 2$ per bit, while
$\langle K \rangle$ is retained as the kinetic energy of the bound state.
The virial theorem $2\langle K \rangle + \langle V \rangle = 0$ is the
mechanical expression of this condition. A stable bound state exists if
and only if both are satisfied simultaneously:}
\begin{equation}
\langle K \rangle = Q = T\,\Delta S = E_\mathrm{Landauer}.
\label{eq:identity_formal}
\end{equation}
\end{quote}

\subsection{Comparison to the virial theorem}

The virial theorem and the thermodynamic identity are compared in
Table~\ref{tab:comparison}.

\begin{table}[H]
\centering
\caption{Comparison of the virial theorem and the thermodynamic identity.}
\label{tab:comparison}
\begin{tabular}{p{3.5cm} p{5.5cm} p{5.5cm}}
\toprule
& \textbf{Virial Theorem} & \textbf{Thermodynamic Identity} \\
\midrule
Statement
& For a stable bound system, $2\langle K\rangle + \langle V\rangle = 0$
& As a system crosses from unbound to bound, $Q = \langle K\rangle =
  T\Delta S = E_\mathrm{Landauer}$; a stable bound state exists
  if and only if both mechanical and thermodynamic conditions are
  satisfied simultaneously \\[6pt]
Derivation
& From Newton's second law (Clausius, 1870)
& From the first law, second law, and Landauer's principle (this work) \\[6pt]
What it describes
& The equilibrium state
& The condition for the equilibrium state to exist \\[6pt]
Physical content of $K$
& Average kinetic energy of bound particles
& Heat released irreversibly at the virialization boundary \\[6pt]
Explains the $1/2$?
& No
& Yes: the unique ratio at which mechanical and thermodynamic
  equilibrium are simultaneously satisfied \\[6pt]
Scale restriction
& None stated; verified mechanically
& None: inherited from the scale-invariance of Landauer's principle \\[6pt]
Predicts non-virialization?
& No: describes only existing bound states
& Yes: systems that cannot encode $Q$ on an available surface
  cannot maintain bound equilibrium \\
\bottomrule
\end{tabular}
\end{table}

\subsection{Why the $1/2$ is no longer unexplained}

The ratio $\langle K \rangle / |\langle V \rangle| = 1/2$ arises because
the $1/r$ potential is homogeneous of degree $-1$. Euler's theorem for
homogeneous functions gives $2\langle K \rangle = |\langle V \rangle|$
mechanically. The thermodynamic identity shows why this mechanical result
has a thermodynamic meaning: the system partitions $|\langle V \rangle|$
exactly in half at equilibrium because that is the unique partition at
which the Landauer cost of the transition equals the kinetic energy of the
resulting bound state. Any other partition would require either net energy
input (unstable) or further energy release (not yet at equilibrium).

\subsection{Scale invariance}

The scale-invariance of the thermodynamic identity follows directly from
the scale-invariance of Landauer's principle. Landauer's principle applies
wherever an irreversible process produces entropy, at any scale, at any
temperature. The virialization boundary exists at every scale where a $1/r$
potential governs a bound system. Therefore the thermodynamic identity
holds at every such scale. This is not an assumption --- it is a
consequence of the universality of the second law of thermodynamics and
Landauer's principle.

% ─────────────────────────────────────────────────────────────────────────────
\section{Evidence Across 37 Orders of Magnitude}
\label{sec:evidence}
% ─────────────────────────────────────────────────────────────────────────────

Tables~\ref{tab:core} and~\ref{tab:extended} present the evidence for the
thermodynamic identity across every physical domain in which bound systems
under $1/r$ potentials have been tested. The range extends from hydrogen
atoms at $10^{-10}$~m to the cosmological coupling constant at $10^{26}$~m
--- 37 orders of magnitude in physical scale.

\begin{table}[H]
\centering
\small
\caption{Core virial theorem domain: the $1/2$ partition verified across
37 orders of magnitude. \colorbox{mcmcrow}{Blue row}: consistent with 17
independent Planck~2018 MCMC chains with zero free parameters.}
\label{tab:core}
\begin{tabular}{p{3.5cm} p{3.0cm} p{1.8cm} p{3.5cm} p{2.8cm}}
\toprule
\textbf{System} & \textbf{Scale / Force} & \textbf{$K/|V|$} &
\textbf{Result} & \textbf{Precision} \\
\midrule
Electron rest mass
& $\lambda_C$, EM
& Fixed point
& $m_e$ to 6.6~ppm
& Zero free parameters \citep{IAM_electron} \\[3pt]
Atoms (H--Xe, 20 elements)
& $10^{-10}$~m, EM
& $T/|V|$
& $1.0000000000$
& Machine precision \citep{Clementi1974,Bunge1993} \\[3pt]
Molecules (H$_2$--CO$_2$, 10)
& $10^{-9}$~m, EM
& $T/|V|$
& $1.0000000000$
& Exact at equil.\ geom.\ \citep{Clementi1974,Bunge1993} \\[3pt]
Equipartition theorem
& $10^{-9}$--$10^{-2}$~m, Thermal
& $k_BT/2$ per DOF
& $1/2$ exact
& Exact in limit (harmonic potential; same ratio, different theorem) \\[3pt]
Solar structure
& $10^{9}$~m, Gravity
& $K/|U|$
& $\approx 0.5$
& $\sim 10\%$ \citep{IAM_virial} \\[3pt]
White dwarfs (Chandrasekhar)
& $10^{7}$~m, Gravity
& $M_\mathrm{Ch}$
& $1.44\,M_\odot$
& Exact, observed \citep{Chandrasekhar1931} \\[3pt]
Galaxy clusters ($\sim$20)
& $10^{23}$~m, Gravity
& $K/|U|$
& $\approx 0.5$
& $\sim 20\%$ \citep{IAM_virial} \\[3pt]
Virial efficiency $\langle\eta\rangle$ (6 N-body)
& $10^{22}$--$10^{25}$~m, Gravity
& $2K/|U|$
& $0.815 \pm 0.025$
& 6 independent sims \citep{BryanNorman1998,Bett2007,Neto2007,Ludlow2010,Power2012,Klypin2016} \\[3pt]
\rowcolor{mcmcrow}
Cosmological coupling $\beta_m$
& $10^{26}$~m, Gravity
& $\beta_m/\Omega_m$
& $0.1583 \pm 0.0033$
& $0.2\sigma$, 17 chains \citep{IAM_obs} \\
\midrule
\textbf{Scale range}
& \multicolumn{4}{l}{$\lambda_C$ (electron) to $10^{26}$~m (cosmic horizon) ---
\textbf{$>$37 orders of magnitude}} \\
\bottomrule
\end{tabular}
\end{table}

\begin{table}[H]
\centering
\small
\caption{Extended domain: expressions of the thermodynamic identity beyond
the classical virial theorem. \colorbox{newresult}{Green rows}: results
derived within the IAM framework. \colorbox{weakrow}{Red row}: structural
exception (see Section~\ref{sec:forces}).}
\label{tab:extended}
\begin{tabular}{p{3.8cm} p{2.5cm} p{4.5cm} p{3.8cm}}
\toprule
\textbf{Domain} & \textbf{Scale} &
\textbf{Expression} & \textbf{Result} \\
\midrule
$H_0$ sector split
& Cosmological
& $H_0^\mathrm{matter}/H_0^\mathrm{photon} = \sqrt{1+\beta_m}$
& 72.26/67.16 km\,s$^{-1}$\,Mpc$^{-1}$, both $<1\sigma$ \citep{IAM_obs} \\[3pt]
$\sigma_8$ suppression
& Cosmological
& $\sigma_8^\mathrm{IAM} = 0.800$
& $0.7998 \pm 0.0058$, $0.1\sigma$ \citep{IAM_obs} \\[3pt]
M--$\sigma$ relation (SMBHs)
& $10^{20}$~m
& $M_\mathrm{BH} \propto \sigma^4$
& 2\% from observed, $0.7\sigma$ (115 galaxies) \citep{IAM_msigma} \\[3pt]
BH entropy partition
& $r_s = 2GM/c^2$
& $E_\mathrm{Landauer} = (\ln 2/2)M_\mathrm{BH}c^2$
& 34.7\% of rest mass (universal) \citep{IAM_BH} \\[3pt]
\rowcolor{newresult}
Dark sector virial partners
& $a = 1$
& $\Omega_m/[\beta_m\mathcal{E}(1)] = 2$ exactly
& Exact, zero free parameters \citep{IAM_virial_dm} \\[3pt]
\rowcolor{weakrow}
\textbf{Weak nuclear force}
& $10^{-18}$~m
& \textbf{Structural exception: cannot satisfy equilibrium condition}
& \textbf{Consistent with identity's physical content (Sec.~\ref{sec:forces})} \\
\bottomrule
\end{tabular}
\end{table}

% ─────────────────────────────────────────────────────────────────────────────
\section{The Four Fundamental Forces}
\label{sec:forces}
% ─────────────────────────────────────────────────────────────────────────────

The thermodynamic identity is satisfied by every force that maintains
stable bound structures. Three of the four fundamental forces fall into
this category.

\textbf{Gravity} operates through a $1/r$ potential in three-dimensional
space. Gauss's law requires the $1/r^2$ force from any point source in 3D,
giving the $1/r$ potential and therefore the exact virial ratio $1/2$ for
any gravitationally bound system. The thermodynamic identity is verified
from stellar structure through galaxy clusters to the cosmological coupling
constant at $0.2\sigma$ from 17 independent MCMC chains.

\textbf{Electromagnetism} operates through the same $1/r$ Coulomb
potential. The thermodynamic identity is verified at machine precision for
20 atoms and 10 molecules from published Hartree--Fock
energies~\citep{Clementi1974,Bunge1993}. The mean virial ratio $\eta = T/|E| =
1.0000000000$ with standard deviation zero. For converged Hartree--Fock
wavefunctions, the virial ratio must equal unity exactly as a mathematical
consequence of the $1/r$ potential --- an exact verification of the
thermodynamic identity at the quantum level.

\textbf{The strong nuclear force} produces the first virialized structures
in cosmic history: protons and neutrons forming at $t \approx 10^{-6}$~s
when the quark-gluon plasma underwent QCD confinement. The resulting
nucleons are bound states satisfying the thermodynamic identity within
the strong-force domain. Note that the Chandrasekhar white dwarf limit
--- though commonly cited in this context --- arises from the equilibrium
between electron degeneracy pressure and gravitational collapse, and is
therefore a gravitational verification of the thermodynamic identity, not
a strong-force one. It is listed in Table~\ref{tab:core} under gravity.

\textbf{The weak nuclear force is the structural exception.} It violates
parity~\citep{Wu1957}, violates CP symmetry~\citep{Cronin1964}, and is
responsible for beta decay, stellar nucleosynthesis, and the
matter-antimatter asymmetry that produced the observable
universe~\citep{Sakharov1967}. It cannot satisfy the thermodynamic
identity not because it fails to do so empirically but because satisfying
it would prevent the weak force from performing its physical function.

The thermodynamic identity describes the equilibrium condition of stable
bound states. The weak force's function is to enable irreversible
transitions between stable states --- to dissolve one equilibrium and
create another. A force that satisfies the equilibrium condition maintains
the current state. The weak force by definition transforms the current
state. These functions are structurally incompatible.

The exception is evidence, not a counterexample. A thermodynamic identity
that holds for all forces governing stable structures and is violated by
precisely the one force whose function is structural transformation
distinguishes equilibrium forces from transition forces on physical
grounds.

% ─────────────────────────────────────────────────────────────────────────────
\section{The MCMC Confirmation}
\label{sec:mcmc}
% ─────────────────────────────────────────────────────────────────────────────

The most rigorous quantitative test of the thermodynamic identity is the
Planck~2018 MCMC posterior returning $\beta_m = 0.1583 \pm 0.0033$ against
the theoretically predicted $\beta_m = \Omega_m/2 = 0.1577$ --- a
confirmation at $0.2\sigma$ --- from 17 independent converged chains with
no free parameters, no discarded runs, and no excluded redshift
ranges~\citep{IAM_obs}.

The identity predicts $\beta_m = \Omega_m/2$ from the virial partition of
gravitational decoherence energy, before any chains were run. The
Planck~2018 MCMC sampler has no knowledge of the thermodynamic identity,
the virial theorem, or the physical argument that $\beta_m = \Omega_m/2$.
It samples the posterior probability distribution of the cosmological
parameters given the Planck~2018 likelihood. The posterior peak at
$\beta_m = 0.1583$ arrives there because the data are consistent with that
value.

This sequence --- derivation from first principles before observation,
the data consistent with the prediction from an independent measurement
with no knowledge of the prediction --- constitutes the standard of
scientific evidence. The
confirmation at $0.2\sigma$ across 17 independent chains is the strongest
available quantitative test of the thermodynamic identity at cosmological
scales.

\begin{table}[H]
\centering
\caption{MCMC confirmation of $\beta_m = \Omega_m/2$. All 17 chains
converged to $R-1 < 0.01$. No chains discarded. Zero free parameters
beyond $\Lambda$CDM.}
\label{tab:mcmc}
\begin{tabular}{p{3.2cm} p{5.2cm} p{5.2cm}}
\toprule
\textbf{Chain set} & \textbf{Dataset} &
\textbf{$\beta_m$ posterior vs prediction ($0.1577$)} \\
\midrule
Level 1 (12 chains)
& Planck~2018 + combinations
& $0.1577$--$0.1592$ ($< 0.3\%$ from prediction) \\[3pt]
Level 2 Run A
& Planck~2018 full likelihood
& $0.1583 \pm 0.0033$ ($0.2\sigma$) \\[3pt]
Level 2 Run D
& Planck~2018 + RSD
& $0.1583 \pm 0.0033$ ($0.2\sigma$) \\
\midrule
\textbf{All 17 chains}
& \textbf{All combinations}
& \textbf{Consistent at $< 0.3\sigma$, zero free parameters} \\
\bottomrule
\end{tabular}
\end{table}

The best-fit $\Delta\chi^2 = +0.54$ relative to $\Lambda$CDM --- with
$\beta_m$ fixed by the thermodynamic identity and zero additional free
parameters --- indicates that the IAM framework, which instantiates this
identity at cosmological scales, is statistically indistinguishable from
$\Lambda$CDM under the full Planck likelihood while simultaneously
providing a physical account of the $\sigma_8$ and $H_0$ tensions that
$\Lambda$CDM does not address~\citep{IAM_obs}.

An independent cross-check from N-body simulations confirms the same
prediction. Six independent simulation
studies~\citep{BryanNorman1998,Bett2007,Neto2007,Ludlow2010,Power2012,Klypin2016}
spanning 25~years of published literature measure the mass-weighted virial
efficiency $\langle\eta_\mathrm{vir}\rangle = 0.815 \pm 0.025$, consistent
with the thermodynamic identity prediction $\eta_\mathrm{vir} =
1/(2f_\mathrm{coll}) = 0.81$ using the published collapsed fraction
$f_\mathrm{coll} = 0.62 \pm 0.03$~\citep{Tinker2008}. The combined
measurement $\Omega_m \times f_\mathrm{coll} \times
\langle\eta_\mathrm{vir}\rangle = 0.315 \times 0.62 \times 0.815 = 0.159
\pm 0.010$ agrees with $\beta_m = \Omega_m/2 = 0.1575$ to $1.0\%$ with
no free parameters~\citep{IAM_virial_eff}.

% ─────────────────────────────────────────────────────────────────────────────
\section{Discussion}
\label{sec:discussion}
% ─────────────────────────────────────────────────────────────────────────────

\subsection{The virial theorem as a special case}

The derivation in Section~\ref{sec:derivation} shows that the virial
theorem is the mechanical expression of a thermodynamic identity. The
identity is more general: it governs the boundary between unbound and bound
states at every scale where a $1/r$ potential produces stable equilibrium.
The virial theorem describes the equilibrium state. The thermodynamic
identity describes the condition for that state to exist and to persist.

This distinction has physical consequences. The virial theorem is silent on
what happens when the conditions for equilibrium cannot be satisfied. The
thermodynamic identity predicts the outcome directly: a system that cannot
encode $Q$ on an available surface cannot maintain bound equilibrium and
returns to the unbound state. This prediction is consistent with the
missing satellites problem, where dark matter halos below a critical
velocity dispersion $\sigma_\mathrm{crit} \approx 4$~km\,s$^{-1}$ disperse
rather than virialize, because no local black hole horizon is available to
absorb the entropy produced at the virialization
boundary~\citep{IAM_satellites}.

\subsection{Landauer's principle as the physical basis for scale invariance}

The virial theorem has been verified across 37 orders of magnitude without
a physical explanation for why it should hold across such a range. The
thermodynamic identity provides that explanation. Landauer's principle
applies wherever an irreversible process produces entropy. The
virialization boundary exists wherever a $1/r$ potential produces a bound
state. The boundary condition $\langle K \rangle = Q$ therefore holds at
every such scale, from the hydrogen atom to the cosmic horizon, because
Landauer's principle is scale-invariant and the $1/r$ potential form is
universal.

This addresses the question of why the virial theorem should apply at
cosmological scales. It is not an extension of a mechanical result beyond
its domain. It is a thermodynamic identity that applies wherever its
conditions are met, at any scale.

\subsection{Relationship to Jacobson 1995}

Jacobson~\citet{Jacobson1995} derived the Einstein equations from
the thermodynamic identity $\delta Q = T\,\delta S$ applied to the Rindler
horizon of an accelerated observer, with $S = A/4\ell_P^2$. That
derivation showed that Einstein's equations are thermodynamic in origin.

The present work applies the same thermodynamic identity
$Q = T\,\Delta S$ to the virialization boundary of a self-gravitating
system and establishes that the virial theorem is its mechanical
expression. Both start from the same thermodynamic identity. Both show
that a fundamental result of classical physics has thermodynamic origin.
The present work extends the Jacobson program from the field equations
governing curved spacetime to the equilibrium condition governing bound
states within that spacetime.

% ─────────────────────────────────────────────────────────────────────────────
\section{Conclusion}
\label{sec:conclusion}
% ─────────────────────────────────────────────────────────────────────────────

The virial theorem $2\langle K\rangle + \langle V\rangle = 0$ has been
verified across 37 orders of magnitude for 154 years without a physical
identification of what $\langle K\rangle$ represents thermodynamically.

This paper shows that the virial theorem and Landauer's principle are two
expressions of the same thermodynamic identity. The first law establishes
$Q = -E_f$; the second law and Landauer's principle establish
$E_\mathrm{Landauer} = Q$; the equilibrium condition of the $1/r$
potential establishes $-E_f = \langle K\rangle$. Together these require:
\begin{equation*}
\langle K\rangle = Q = T\,\Delta S = E_\mathrm{Landauer}.
\end{equation*}
$\langle K\rangle$ is the heat released irreversibly at the virialization
boundary, encoded as entropy on the nearest available surface at Landauer
cost $k_BT\ln 2$ per bit. The virial theorem is the mechanical expression
of this thermodynamic identity.

The virial theorem is the mechanical expression of this thermodynamic
identity. The $1/2$ partition is the unique ratio at which mechanical and
thermodynamic equilibrium are simultaneously satisfied for a $1/r$
potential. Its appearance across 37 orders of magnitude is a consequence
of the scale-invariance of Landauer's principle and the universality of
the $1/r$ potential form.

Three of the four fundamental forces satisfy the identity. The weak nuclear
force is the structural exception, excluded by physical necessity. The
data are consistent with the identity at $0.2\sigma$ at cosmological
scales from 17 independent MCMC chains with zero free parameters.

We propose that this thermodynamic identity constitutes a fundamental law
of physics of which the virial theorem is the mechanical special case.

% ─────────────────────────────────────────────────────────────────────────────
\section*{Acknowledgments}
% ─────────────────────────────────────────────────────────────────────────────

The thermodynamic framework of T.\ Jacobson~\citep{Jacobson1995} is the
foundation upon which this work is built. The author thanks
R.\ Landauer~\citep{Landauer1961} and R.\ Clausius~\citep{Clausius1870}
for the tools that make this connection possible.

% ─────────────────────────────────────────────────────────────────────────────
\section*{Data Availability}
% ─────────────────────────────────────────────────────────────────────────────

All MCMC chains, analysis scripts, and supporting papers are publicly
archived at \url{https://github.com/hmahaffeyges/IAM-Validation} with
canonical Zenodo DOI
\href{https://doi.org/10.5281/zenodo.18702042}{10.5281/zenodo.18702042}.

% ─────────────────────────────────────────────────────────────────────────────
\bibliographystyle{plainnat}
\bibliography{iam_thermodynamic_identity}
% ─────────────────────────────────────────────────────────────────────────────

\end{document}
